\section{Introduction}

Recent advances in large language models (LLMs) have demonstrated remarkable capabilities in financial analysis, from sentiment extraction to price forecasting \cite{lopez2023can,chen2023fingpt}. However, a fundamental question remains: Can these models genuinely understand market microstructure mechanisms, or are they sophisticated pattern-matching systems that recall memorized associations from training data?

This distinction matters. Market makers face regulatory requirements to maintain delta-neutral positions despite accumulating gamma exposure \cite{sec15c31}. When net gamma exposure becomes significantly negative, dealers must hedge pro-cyclically—selling into rallies and buying into dips—amplifying price movements through forced trading flows \cite{ni2005stock,garleanu2009demand}. Testing whether LLMs understand these structural constraints presents a unique challenge: if we present an LLM with "SPY on January 27, 2021" and substantial negative gamma exposure, any successful prediction might simply reflect memorization of the GameStop squeeze rather than reasoning about underlying mechanics. This training data leakage problem pervades current LLM evaluation in finance, where temporal context enables recall rather than reasoning.

We address this through \textit{obfuscation testing}, a novel validation methodology that strips all memorizable context while preserving structural relationships. By converting dates to generic labels ("Day T+0"), removing ticker identities ("SPY" becomes "INDEX\_1"), and eliminating economic event references, we force the LLM to reason solely from market structure—gamma exposure levels, strike distributions, and open interest concentrations. If an LLM can still detect dealer hedging patterns under these conditions, it demonstrates understanding of causal mechanisms rather than pattern matching from training data.

Our approach makes four key contributions. First, we introduce obfuscation testing for validating LLM structural reasoning in financial markets, distinguishing genuine understanding from memorization. Second, we quantify the impact of prompt bias on pattern detection: including regime labels (e.g., "NEGATIVE\_GAMMA") inflates detection rates from 71.5\% to 100\%, revealing how contextual hints create inflated performance metrics. Third, we validate our methodology across multiple pattern types and time periods, demonstrating generalization rather than overfitting. Fourth, we establish that detection capability remains stable while economic profitability varies, proving LLMs identify structural constraints independent of trading value.

Testing this framework on 242 trading days throughout 2024, we evaluate three manifestations of dealer hedging constraints: gamma positioning, stock pinning, and 0DTE hedging. Using fully unbiased prompts with obfuscated data, we achieve 71.5\% average detection rate, significantly exceeding the 60\% mechanical threshold. More importantly, 91.2\% of detected patterns materialize in forward returns, validating that the LLM identifies genuine market mechanics rather than spurious correlations.

These findings have important implications for both AI validation and financial market analysis. For the AI community, obfuscation testing offers a rigorous framework for validating reasoning capabilities in domains where training data contamination is a concern. For finance practitioners, our results suggest LLMs can serve as sophisticated pattern recognition tools for market microstructure analysis, provided their limitations are understood.

The remainder of this paper is organized as follows. Section II reviews related work. Section III presents our methodology. Section IV describes the experimental setup. Section V reports results. Section VI discusses implications and limitations. Section VII concludes.